% Socle Identitaire V13.1-FR — Version académique française
% Version originale en anglais : identity_bedrockV13.1.tex
% Compilation : pdflatex -> pdflatex

\documentclass[11pt,a4paper]{article}
\usepackage[utf8]{inputenc}
\usepackage[T1]{fontenc}
\usepackage{lmodern}
\usepackage{microtype}
\usepackage[french]{babel}
\usepackage{amsmath,amssymb,amsthm}
\usepackage{mathtools}
\usepackage{booktabs}
\usepackage{enumitem}
\usepackage[colorlinks=true,linkcolor=black,citecolor=black,urlcolor=blue]{hyperref}
\usepackage{geometry}
\geometry{a4paper,margin=25mm}

\newtheorem{theoreme}{Théorème}[section]
\newtheorem{definition}[theoreme]{Définition}
\newtheorem{corollaire}[theoreme]{Corollaire}
\renewcommand{\proofname}{Démonstration}

\title{\textbf{Socle Identitaire : Un Substrat Cryptographique Minimal\\
et ses Implications Systémiques}}
\author{Jean-François Bouchard\\
  \small PrivateNess Network (co-fondateur, Directeur technique)\\
  \small Emercoin (contributeur communautaire)\\
  \small \texttt{[contact via privateness.network]}}
\date{Mars 2026}

\begin{document}
\maketitle

\noindent\textbf{Divulgation relative à l'IA.} L'auteur a eu recours à des
outils de traitement automatique du langage pour la rédaction et la révision
du manuscrit.  Les contributions techniques, la conception du système, les
réalisations, le cadrage théorique et les idées directrices décrits ci-après
sont issus de l'auteur, qui s'appuie sur neuf années d'expérience
opérationnelle dans le déploiement d'infrastructures fondées sur Emercoin.

\begin{abstract}
Nous présentons le Socle Identitaire, un système d'identité déployé construit
à partir d'un substrat cryptographique minimal.  En partant de l'exigence
qu'aucune clé institutionnelle n'existe sous la clé individuelle, nous
dérivons une architecture fondée sur les fonctions à sens unique (FSU), qui
sont nécessaires et suffisantes pour les schémas de signature à
infalsifiabilité existentielle sous attaque adaptative à messages choisis.

Le système sépare la cryptographie d'identité du consensus de répertoire au
moyen d'une couche d'ancrage en ajout exclusif, indépendante de la chaîne de
blocs, permettant la portabilité entre substrats hétérogènes tout en
préservant la vérifiabilité.

Nous montrons en outre que des capacités système supplémentaires---nommage,
représentation structurée, récupération de contenu, exécution déterministe,
connectivité et interaction économique---émergent comme propriétés dérivées de
ce substrat identitaire plutôt que comme sous-systèmes indépendants.

La contribution n'est pas un nouveau primitif, mais une composition de
primitifs existants en un système qui réduit les hypothèses de confiance au
minimum irréductible tout en permettant un déploiement pratique sur
infrastructures matérielles grand public.
\end{abstract}

\noindent\textbf{Mots-clés :} fonctions à sens unique $\cdot$ signatures
numériques $\cdot$ identité souveraine $\cdot$ annuaire en ajout exclusif
$\cdot$ propriétés émergentes $\cdot$ architecture indépendante de la chaîne
de blocs $\cdot$ déploiement opérationnel

%--------------------------------------------------------------------
\section{Introduction}
%--------------------------------------------------------------------

Les systèmes d'identité numérique font communément appel à des ancres de
confiance institutionnelles---autorités de certification, fournisseurs
d'identité, opérateurs de registre.  Ces entités introduisent des points de
contrôle extérieurs sur l'identité et créent des dépendances affectant la
disponibilité, la portabilité et la souveraineté de l'utilisateur.

Ce travail examine si l'identité peut être construite en réduisant au minimum
de telles dépendances, en ne reposant que sur des hypothèses cryptographiques
nécessaires.  Nous examinons également les implications plus larges d'une
telle construction, en montrant que de multiples capacités au niveau système
peuvent être dérivées d'un substrat identitaire minimal.

\subsection{Contributions}

\begin{itemize}[leftmargin=*]
  \item \textbf{Construction minimale.}  Identité fondée sur des schémas de
    signature ancrés dans les FSU.
  \item \textbf{Architecture indépendante de la chaîne de blocs.}  Séparation
    de la cryptographie d'identité du consensus de répertoire.
  \item \textbf{Système déployé.}  Mise en œuvre publiquement vérifiable,
    opérationnelle depuis 2020 sur le réseau PrivateNess~\cite{onchain}.
  \item \textbf{Émergence systémique.}  Démonstration que les capacités de
    couche supérieure émergent du substrat identitaire sans racines de
    confiance supplémentaires.
\end{itemize}

%--------------------------------------------------------------------
\section{Préliminaires}
%--------------------------------------------------------------------

Les fonctions à sens unique sont nécessaires et suffisantes pour les schémas
de signature à infalsifiabilité existentielle sous attaque adaptative à
messages choisis~\cite{rompel1990}.  Cette équivalence contraint toute
construction reposant sur ce primitif à une unique hypothèse d'intractabilité,
et aucune hypothèse strictement plus faible n'est suffisante.

Les systèmes d'identité traditionnels outrepassent cette exigence minimale en
introduisant des modèles de confiance hiérarchiques.  L'infrastructure à clé
publique (ICP)~\cite{rfc5280} en est l'exemple canonique~: les autorités de
certification peuvent délivrer des certificats erronés, être compromises ou
contraintes, introduisant des modes de défaillance absents du primitif
sous-jacent.  Les systèmes d'identité fédérés (OAuth~2.0, OpenID~Connect)
délèguent l'identité à un fournisseur tiers susceptible de désauthentifier
les utilisateurs unilatéralement.

\begin{definition}[Schéma de signature EUF-AAMC]
  \label{def:sig}
  Un \emph{schéma de signature} $\Sigma = (\mathsf{KeyGen}, \mathsf{Sign},
  \mathsf{Verify})$ est à \emph{infalsifiabilité existentielle sous attaque
  adaptative à messages choisis} (EUF-AAMC) si pour tout adversaire PPT
  $\mathcal{A}$ ayant accès à un oracle de signature,
  \[
    \Pr[\mathcal{A}\ \text{produit une falsification valide}]
    \leq \mathrm{negl}(\lambda),
  \]
  où $\lambda \in \mathbb{N}$ est le paramètre de sécurité et $\mathrm{negl}$
  une fonction négligeable.
\end{definition}

\begin{definition}[Annuaire en ajout exclusif]
  \label{def:dir}
  Un \emph{annuaire en ajout exclusif} $\mathcal{D}$ est un ensemble de
  paires $(id, pk)$ muni d'un prédicat de liaison $\mathsf{Lier}(id, pk)$
  tel que toute paire une fois inscrite ne peut plus être supprimée par
  aucune procédure efficace.  $\mathcal{D}$ possède la \emph{finalité} si
  la probabilité de réversion d'une entrée à la profondeur $d$ est négligeable
  en $\lambda$ et décroissante en $d$.
\end{definition}

%--------------------------------------------------------------------
\section{Architecture}
%--------------------------------------------------------------------

\subsection{Couche d'identité}

Un participant génère une paire de clés
$(sk, pk) \gets \mathsf{KeyGen}(1^\lambda)$ à partir d'une graine à haute
entropie.  L'authentification consiste à signer un défi et à vérifier la
signature à l'aide de la clé publique~:

\begin{enumerate}[leftmargin=*]
  \item Le vérificateur émet un défi $c \gets \{0,1\}^\lambda$.
  \item Le prouveur calcule $\sigma \gets \mathsf{Sign}(sk, c)$ et renvoie $\sigma$.
  \item Le vérificateur accepte si et seulement si $\mathsf{Verify}(pk, c, \sigma) = 1$.
\end{enumerate}

Aucun certificat, jeton de session ou oracle tiers n'intervient dans cet échange.

\subsection{Substrat de répertoire}

Un annuaire en ajout exclusif lie les identifiants aux clés publiques.  Les
exigences sont les suivantes~:

\begin{itemize}[leftmargin=*]
  \item \textbf{Vérifiabilité globale.}  Tout nœud peut vérifier
    l'appartenance indépendamment.
  \item \textbf{Persistance.}  Les liaisons publiées ne peuvent être supprimées.
  \item \textbf{Finalité.}  Une liaison confirmée à profondeur suffisante est
    irréversible avec probabilité négligeable en $\lambda$.
\end{itemize}

La couche d'identité et la couche de répertoire sont découplées~: tout
substrat satisfaisant ces exigences constitue une instanciation valide,
permettant la portabilité entre substrats hétérogènes sans modification de
la couche cryptographique.

%--------------------------------------------------------------------
\section{Analyse de sécurité}
%--------------------------------------------------------------------

\begin{theoreme}[Solidité de l'authentification]
  \label{thm:solidite}
  Si le schéma de signature $\Sigma$ satisfait la Définition~\ref{def:sig}
  et si l'annuaire $\mathcal{D}$ satisfait la Définition~\ref{def:dir},
  alors l'authentification est infalsifiable sous les hypothèses standard.
\end{theoreme}

\begin{proof}
Une falsification réussie implique soit (i) une rupture du schéma de
signature---produire une signature valide sans détenir $sk$---soit (ii) une
violation de l'intégrité de l'annuaire---substituer $(id, pk')$ à $(id, pk)$
afin de tromper le vérificateur sur la clé publique légitime.  Le cas (i)
contredit la sécurité EUF-AAMC de $\Sigma$~; le cas (ii) contredit la
propriété d'ajout exclusif avec finalité de $\mathcal{D}$.  Les deux cas
étant exclus par hypothèse, aucune falsification n'est réalisable.
\end{proof}

\begin{corollaire}
  Toute usurpation réussie de l'identité $id$ exige soit de casser le
  schéma de signature, soit de corrompre l'annuaire.  Aucune de ces
  conditions ne fait intervenir un secret institutionnel.
\end{corollaire}

%--------------------------------------------------------------------
\section{Propriétés système émergentes}
%--------------------------------------------------------------------

Dès lors qu'un substrat identitaire satisfaisant les
Définitions~\ref{def:sig} et~\ref{def:dir} existe, des capacités système
supplémentaires deviennent disponibles comme propriétés dérivées, sans
introduction de nouvelles racines de confiance.

\subsection{Nommage et résolution}

L'identité permet un nommage décentralisé via des liaisons en ajout exclusif
entre identifiants intelligibles et clés publiques.  La résolution devient un
processus de vérification de l'appartenance plutôt que de confiance envers
des autorités hiérarchiques.  Pour résoudre le nom $n$, il suffit de
récupérer $(n, pk)$ dans $\mathcal{D}$ et de vérifier $\mathsf{Lier}(n, pk) = 1$.

\subsection{Représentation structurée des données}

Des objets peuvent être représentés selon des formats structurés déterministes.
La sérialisation canonique permet l'analyse, l'indexation et l'interrogation
cohérentes entre nœuds indépendants, rendant possible la découverte
décentralisée sans autorité de schéma partagée.

\subsection{Adressage par contenu}

La récupération de contenu est effectuée à l'aide de hachés cryptographiques~:
pour une donnée $m$, son adresse est $h = H(m)$ et tout récupérateur vérifie
$H(m') = h$ avant utilisation.  L'intégrité est ainsi indépendante du
mécanisme de transport~\cite{ipfs}.

\subsection{Exécution déterministe}

L'exécution peut être modélisée comme une fonction pure~:
\[
  h_{\mathrm{sortie}} = H\bigl(F(m_1, \ldots, m_k)\bigr),
\]
où $F$ est un programme adressé par contenu et $\{m_i\}$ ses entrées,
elles-mêmes adressées par contenu.  Cela permet un calcul reproductible et
des résultats vérifiables indépendamment par tout nœud ayant accès aux entrées.

\subsection{Connectivité en réseau superposé}

Les techniques de réseau superposé permettent aux nœuds de communiquer
indépendamment des contraintes du réseau sous-jacent.  L'acheminement fondé
sur l'identité réduit la dépendance à l'égard d'une infrastructure de
routage centralisée~; chaque point d'extrémité étant authentifié via le
substrat d'identité, le réseau superposé hérite de la sécurité de la couche
d'identité sans ICP séparée.

\subsection{Interaction économique}

Des mécanismes d'incitation peuvent être construits à l'aide de l'attribution
liée à l'identité, permettant des modèles de récompense fondés sur la
participation sans coordination centralisée.  Chaque action économique est
signée par la clé secrète de l'acteur, rendant les contributions
attribuables et auditables sans opérateur de registre centralisé.

%--------------------------------------------------------------------
\section{Démonstration pratique}
%--------------------------------------------------------------------

L'architecture a été mise en œuvre et déployée.  Le système est opérationnel
depuis 2020 sur le réseau PrivateNess~\cite{onchain}.

\subsection{Déploiement applicatif}

Les applications peuvent être résolues via le nommage décentralisé,
récupérées par adressage par contenu, et exécutées dans un environnement
déterministe.  L'ensemble du flux ne requiert aucun intermédiaire de confiance.

\subsection{Déploiement nodal}

Les nœuds fonctionnent sur du matériel grand public, incluant des dispositifs
à faible coût (classe RK3528) et des systèmes réaffectés.  Un coût de
déploiement complet vérifié inférieur à 50~CAD a été constaté~\cite{onchain},
validant la faisabilité sur des infrastructures sans spécialisation.

\subsection{Flux de résolution}

Une interaction typique suit les étapes~:

\begin{enumerate}[leftmargin=*]
  \item Résoudre l'identifiant en liaison clé publique dans l'annuaire.
  \item Récupérer le contenu associé par son haché.
  \item Vérifier l'intégrité~: $H(m') = h$.
  \item Exécuter de façon déterministe.
\end{enumerate}

Chaque étape est indépendamment vérifiable par tout nœud disposant d'un
accès à l'annuaire public et aux fonctions de hachage standard.

%--------------------------------------------------------------------
\section{Mise en œuvre}
%--------------------------------------------------------------------

Le système est mis en œuvre à l'aide d'un substrat nom--valeur en ajout
exclusif ancré à des mécanismes de consensus externes.  Les liaisons
d'identité sont publiquement vérifiables et portables entre substrats.

Concrètement, l'annuaire public $\mathcal{D}$ est instancié à l'aide du
Stockage Nom--Valeur (SNV) d'Emercoin~\cite{emercoin}, un magasin
clé--valeur natif de la chaîne de blocs Emercoin, ancré à Bitcoin via
Preuve de Travail Auxiliaire (AuxPoW)~\cite{hack51}.  Les enregistrements
d'identité sont écrits sous des clés SNV de la forme
\texttt{id:$\langle$identifiant$\rangle$} selon le format WORM (\emph{Write-Once
Read-Many})~\cite{worm}.  Une fois écrits, les enregistrements peuvent être
renouvelés mais non supprimés, satisfaisant la Définition~\ref{def:dir}.

%--------------------------------------------------------------------
\section{Évaluation}
%--------------------------------------------------------------------

L'authentification n'exige que la vérification de signature et une
consultation de l'annuaire, ce qui se traduit par un faible coût calculatoire
et une latence minimale.

\textbf{Complexité calculatoire.}  La vérification d'une authentification
exige une vérification de signature Ed25519~\cite{eddsa} ($O(1)$ en taille
du message) et une consultation de l'annuaire SNV ($O(\log n)$ dans l'arbre
de Merkle).  Ces deux opérations sont réalisables sur du matériel embarqué
de classe RK3528.

\textbf{Stabilité opérationnelle.}  L'enregistrement d'identité canonique
en chaîne~\cite{onchain} est ancré dans le bloc SNV Emercoin \#445\,669
(1\textsuperscript{er}~novembre~2020, 22:16:45~UTC), transitivement ancré à
Bitcoin via AuxPoW, et auditable à l'adresse
\url{https://explorer.emercoin.com/dpo/PrivateNESS.Network}.  Cinq années
d'exploitation continue (2020--2026) sans interruption de service liée à la
couche cryptographique valident l'adéquation du substrat minimal pour des
déploiements en production.

%--------------------------------------------------------------------
\section{Discussion}
%--------------------------------------------------------------------

L'architecture n'introduit pas de sous-systèmes indépendants pour le nommage,
l'exécution ou la connectivité.  Ces capacités émergent au contraire d'un
substrat identitaire unifié, réduisant la complexité conceptuelle tout en
préservant la compatibilité avec les protocoles existants.

\textbf{Limites du périmètre.}  Le modèle traite des adversaires
calculatoires classiques.  La compromission du point d'extrémité, l'ingénierie
sociale et les attaques quantiques sur la couche de signature sont hors
périmètre.  La migration post-quantique vers ML-DSA-65 (FIPS~204~\cite{fips204})
est envisageable sans modifier les couches de répertoire ou de connectivité.

\textbf{Généralité de la couche de répertoire.}  La Définition~\ref{def:dir}
étant intentionnellement abstraite, toute implémentation satisfaisant les
propriétés d'ajout exclusif et de finalité constitue une instanciation
valide---chaînes de blocs alternatives, journaux distribués vérifiables, ou
systèmes de fichiers en ajout exclusif pour des cas d'usage à moindre risque.

%--------------------------------------------------------------------
\section{Conclusion}
%--------------------------------------------------------------------

Le Socle Identitaire démontre qu'un substrat cryptographique minimal---un
schéma de signature EUF-AAMC et un annuaire en ajout exclusif avec
finalité---est suffisant non seulement pour l'authentification, mais aussi
pour permettre un ensemble complet de capacités système sans introduction
d'ancres de confiance institutionnelles.

En dérivant le comportement système d'un primitif unifié, l'approche réduit
les hypothèses de confiance au minimum tout en permettant un déploiement
pratique et l'interopérabilité avec les protocoles existants.

%--------------------------------------------------------------------
\section*{Remerciements}
%--------------------------------------------------------------------

L'auteur remercie Oleg Khovayko pour les travaux fondateurs sur
l'infrastructure Emercoin et le cadre hack51~\cite{hack51}~; Aleksej Sokolov
(Alex014) pour la conception du format WORM~\cite{worm} et la mise en œuvre
de la chaîne d'outils PrivateNess~; Brandon \og~Synth~\fg~Smietana et Moses
Narrow (\texttt{moe\_narrow} / \texttt{0pcom}) pour le transport Skywire et
l'infrastructure de consensus Fiber~; et la communauté PrivateNess Network
depuis la fondation publique du projet le
1\textsuperscript{er}~septembre~2020~\cite{bitcointalk}.

%--------------------------------------------------------------------
\begin{thebibliography}{99}

\bibitem{rompel1990}
  J.~Rompel,
  \og~One-way functions are necessary and sufficient for secure signatures,~\fg
  in \textit{Proc.\ 22\textsuperscript{e} Symposium annuel ACM STOC},
  pp.~387--394, 1990.

\bibitem{goldreich2001}
  O.~Goldreich,
  \textit{Foundations of Cryptography, Volume~1 : Basic Tools},
  Cambridge University Press, 2001.

\bibitem{eddsa}
  S.~Josefsson et I.~Liusvaara,
  \og~Edwards-Curve Digital Signature Algorithm (EdDSA),~\fg
  RFC~8032, IETF, janvier 2017.
  \url{https://www.rfc-editor.org/rfc/rfc8032}

\bibitem{fips204}
  National Institute of Standards and Technology,
  \og~Module-Lattice-Based Digital Signature Standard (ML-DSA),~\fg
  FIPS~204, août 2024.
  \url{https://doi.org/10.6028/NIST.FIPS.204}

\bibitem{rfc5280}
  D.~Cooper et al.,
  \og~Internet X.509 Public Key Infrastructure Certificate and CRL Profile,~\fg
  IETF RFC~5280, mai 2008.

\bibitem{emercoin}
  Équipe de développement Emercoin,
  \og~Emercoin~: chaîne de blocs hybride avec minage fusionné,~\fg
  \url{https://emercoin.com}, 2013--2026.

\bibitem{hack51}
  O.~Khovayko,
  \og~Estimation du coût d'une attaque à 51\,\% sur Emercoin~: le coefficient
  hack51~\fg [en russe],
  Dzen / Yandex, 2021.
  \url{https://dzen.ru/a/YUotZkN9nAeElPEQ}

\bibitem{worm}
  A.~Sokolov (Alex014) / NESS-Network,
  \og~WORM --- World Object MaRkup language,~\fg
  \url{https://github.com/ness-network/worm}

\bibitem{onchain}
  PrivateNess Network,
  \og~PrivateNess.Network --- enregistrement d'identité canonique en chaîne,~\fg
  Emercoin NVS DPO, bloc~\#445\,669,
  1\textsuperscript{er}~novembre~2020, 22:16:45~UTC.
  \url{https://explorer.emercoin.com/dpo/PrivateNESS.Network}

\bibitem{bitcointalk}
  PrivateNess Network,
  \og~PrivateNess Network (NESS) Privacy CX Fiber Coin,~\fg
  Bitcointalk Announcements, 1\textsuperscript{er}~septembre~2020.
  \url{https://bitcointalk.org/index.php?topic=5272741.0}

\bibitem{eip4361}
  B.~Talsma et al.,
  \og~EIP-4361~: Sign-In with Ethereum,~\fg
  Proposition d'amélioration Ethereum, octobre 2021.
  \url{https://eips.ethereum.org/EIPS/eip-4361}

\bibitem{ipfs}
  J.~Benet,
  \og~IPFS---Content Addressed, Versioned, P2P File System,~\fg
  arXiv:1407.3561, 2014.

\end{thebibliography}

\end{document}
